\documentclass[11pt]{article}
\usepackage[T1]{fontenc}
\usepackage{lmodern,amsmath,amssymb}
\usepackage[margin=1in]{geometry}
\usepackage[hidelinks]{hyperref}
\begin{document}

\title{A Conditional Three-Prime Kolyvagin-Class Test
for the Elliptic Curve 234446.a1}
\author{\small \href{mailto:admin@cosmosgeek.org}{admin@cosmosgeek.org}\qquad \href{mailto:b321014@iiit-bh.ac.in}{b321014@iiit-bh.ac.in}}
\date{Computational working paper}
\maketitle

\begin{abstract}
We record an explicit mod-$7$ localization calculation
for the elliptic curve 234446.a1 over
$K=\mathbb Q(\sqrt{-11})$. Three Kolyvagin primes
select a one-dimensional rational Kummer line.
Under explicit Selmer and lower-support hypotheses,
the corresponding derived Heegner class belongs to that line.
Its nonvanishing is equivalent to a finite-field test at 19.
The test's final scalar remains unevaluated.
\end{abstract}

\section{Curve, points, and coordinates}

Let
\[
 E:\quad y^2+xy=x^3-x^2-79x+289,
 \qquad N_E=234446=2\cdot117223.
\]
Use the integral short Weierstrass coordinates
\[
 X=4x-1,\qquad Y=8y+4x,
\]
giving
\[
 E':\quad Y^2=X^3-1267X+17230.
\]

The conductor and global rank four are recorded in \cite{LMFDB}.
Fix the ordered rational points
\[
\begin{aligned}
 P_1&=(0,17),&
 P_2&=(1,14),\\
 P_3&=(3,7),&
 P_4&=(-4,25).
\end{aligned}
\]
Their coordinates on $E'$ are respectively
\[
 (-1,136),\quad(3,116),\quad
 (11,68),\quad(-17,184).
\]

The group law gives the exact identity
\[
 \boxed{
 B=P_1+P_3=(52/9,-95/27).
 }
\]
In short Weierstrass coordinates,
\[
 B'=(199/9,-136/27).
\]

The points are independent modulo
$7E(\mathbb Q)$, as certified by the localization
matrix below. Under the four-dimensional
plus-Selmer hypothesis, their Kummer images
form a basis of the plus-Selmer space.

\section{Finite-field Kummer calculations}

Put $p=7$. For each prime $\ell$ in the table,
write
\[
 h_\ell=\#E(\mathbb F_\ell)/7.
\]
Choose the generator
\[
 v_\ell=[h_\ell]\overline{P}_{j_\ell}
 \in E(\mathbb F_\ell)[7],
\]
where $j_{1091}=2$ and $j_\ell=1$ otherwise.
All these generators are nonzero.

Define $m_{\ell i}\in\mathbb F_7$ by
\[
 [h_\ell]\overline{P}_i=[m_{\ell i}]v_\ell.
\]

For these primes the rational $7$-primary subgroup is cyclic. Indeed
Frobenius has determinant $\ell\bmod 7\ne1$ on $E[7]$, so cannot fix
all of $E[7]$; each displayed group has order divisible by seven.
Consequently the kernel of the map $R\mapsto[h_\ell]R$ is exactly
$7E(\mathbb F_\ell)$. Direct finite-field group arithmetic gives
\[
\begin{array}{c|r|r|rrrr}
 \ell&a_\ell&\#E(\mathbb F_\ell)&
 m_{\ell1}&m_{\ell2}&m_{\ell3}&m_{\ell4}\\ \hline
 19&-8&28&1&1&5&3\\
 167&14&154&1&2&3&1\\
 853&-28&882&1&2&6&4\\
 1091&-42&1134&0&1&0&3\\
 1553&21&1533&1&3&6&1
\end{array}
\]

For the three proposed support primes and the
original detector, in that order, the matrix is
\[
 M=
 \begin{pmatrix}
 1&2&6&4\\
 0&1&0&3\\
 1&3&6&1\\
 1&2&3&1
 \end{pmatrix},
 \qquad
 \det M=-18\equiv3\pmod7.
\]

The first three rows have rank three, with
\[
 \ker M_{\{853,1091,1553\}}
   =\mathbb F_7(1,0,1,0).
\]
To reproduce the table, count
$1+\sum_{X\bmod\ell}(1+(\frac{X^3-1267X+17230}{\ell}))$
points, multiply by $h_\ell$ using the elliptic group law, and compare
the result with the seven multiples of $v_\ell$. In order of the
five table rows, the generators on the short model are
$(2,6),(166,136),(537,642),(967,954),(1226,130)$.
Therefore the rational Kummer classes with
zero localization at all three support primes
form precisely
\[
 \boxed{\mathbb F_7\,\delta_7(B)}
\]
inside the four-dimensional rational Kummer
space.

The remaining detector coordinates are
\[
 m_{167}(B)=4,\qquad m_{19}(B)=6.
\]
Thus both 167 and 19 detect every nonzero class
on this line.

The comparison with the local Kummer spaces
over $K$ is valid because the listed support
primes are inert in $K$, and extension of finite
residue fields from degree one to degree two
does not destroy mod-$7$ divisibility
information on rational residue points.

\section{Ring-class group and derivative}

Let
\[
 K=\mathbb Q(\sqrt{-11}),\qquad
 n_*=853\cdot1091\cdot1553
     =1445257519.
\]
The field $K$ has class number one and unit group
$\{\pm1\}$. Each prime dividing $n_*$ is inert
in $K$.

Enumerating odd primes $\ell\le1553$ with
$\ell\equiv-1\bmod7$, $(\frac{-11}{\ell})=-1$ and $7\mid a_\ell$
gives exactly $167,853,1091,1553$. Trial division and the point-count
formula above suffice for the enumeration. Since $167$ detects $B$,
the last three are the smallest three admissible support primes at
which $B$ is locally divisible by seven. This proves minimality among
squarefree products of three primes satisfying these requirements.

Let $L=K[n_*]$ be the ring-class field of
conductor $n_*$. The ring-class number formula
gives
\[
 G=\operatorname{Gal}(L/K)
 \simeq C_{854}\times C_{1092}\times C_{1554},
\]
and
\[
 \boxed{[L:K]=1449210672.}
\]

The unit index in the ring-class formula
$[\mathcal O_K^\times:\mathcal O_{n_*}^\times]$
equals one.

For each $\ell\mid n_*$, choose a generator
$\sigma_\ell$ of the local cyclic factor
\[
 \frac{(\mathcal O_K/\ell\mathcal O_K)^\times}
       {(\mathbb Z/\ell\mathbb Z)^\times}
 \simeq C_{\ell+1},
\]
using a fixed Artin-reciprocity convention.

Write $\ell+1=7d_\ell$, where
\[
 (d_{853},d_{1091},d_{1553})
       =(122,156,222).
\]
Let
\[
 H=\prod_{\ell\mid n_*}\langle\sigma_\ell^7\rangle.
\]
Then
\[
 |H|=4225104\equiv2\pmod7,\qquad
 G/H\simeq C_7^3.
\]

Let $g_\ell$ denote the image of $\sigma_\ell$
in $G/H$, and put $T_\ell=g_\ell-1$.

For $M_\ell=\ell+1$, define
\[
 D_\ell=\sum_{j=1}^{M_\ell-1}
             j\sigma_\ell^j,\qquad
 N_\ell=\sum_{j=0}^{M_\ell-1}\sigma_\ell^j.
\]
These satisfy
\[
 (\sigma_\ell-1)D_\ell=M_\ell-N_\ell.
\]

Modulo seven,
\[
 D_\ell=
 \left(\sum_{r=1}^{6}r\sigma_\ell^r\right)
 \left(\sum_{t=0}^{d_\ell-1}
                      \sigma_\ell^{7t}\right).
\]

In $\mathbb F_7[C_7]$ the seven-term derivative
has the exact expression
\[
 D_7(g)=\sum_{r=1}^{6}rg^r
       =-g(g-1)^5.
\]

The image of the full derivative
$D_{n_*}=\prod_{\ell\mid n_*}D_\ell$
in the quotient group algebra is
\[
 \overline D_{n_*}
   =5g_{853}g_{1091}g_{1553}
     T_{853}^5T_{1091}^5T_{1553}^5.
\]
Here the scalar $5$ comes from
$d_{853}d_{1091}d_{1553}\equiv2$
and the three minus signs.

\section{Compression of the Heegner point}

Let $y_{n_*}\in E(L)$ be the Heegner point
obtained from the specified Shimura-curve
parametrization and its fixed CM orientation.

Define the unnormalized trace
\[
 Y_H=\operatorname{Tr}_H(y_{n_*})
       =\sum_{h\in H}h\,y_{n_*}.
\]
It belongs to $E(L^H)$.

The derived point
\[
 z=D_{n_*}y_{n_*}
\]
has the following exact expression modulo
$7E(L)$:
\[
 \boxed{
 z\equiv
 \sum_{r_1,r_2,r_3=1}^{6}
 r_1r_2r_3\,
 g_1^{r_1}g_2^{r_2}g_3^{r_3}Y_H
 \pmod{7E(L)}.
 }
\]
There are $6^3=216$ nonzero terms.

No division by $[L:K]$ or $|H|$ is used.
If the normalized average $|H|^{-1}Y_H$
is used instead, its factor
$|H|\equiv2\pmod7$ must be restored.

The Heegner norm relations and the identities
above show that the class of $z$ in $E(L)/7$
is $G$-invariant.

Under the full residual-image hypothesis, the image of $G_K$
contains $\operatorname{SL}_2(\mathbb F_7)$: the quotient of
the full image by the image of $G_K$ is abelian, since $K/\mathbb Q$
is quadratic. The latter special linear group is perfect. As $L/K$
is abelian, the image of $G_L$ still contains this group.
Thus $E[7](L)=0$. Inflation--restriction
therefore gives
\[
 H^1(K,E[7])
   \simeq H^1(L,E[7])^G.
\]
Hence the invariant point class has a unique
global cohomological descent, without
dividing by the ring-class degree.

Identifying this descent with the standard
mod-seven Kolyvagin class requires the
specified Heegner normalization and local
finite--singular comparison.

\section{Conditional rational-point line}

Let $\kappa_n\in H^1(K,E[7])$ denote the mod-seven system for a
fixed integral Heegner parametrization. For this section assume that
its local conditions and finite--singular recursion are those in
\cite[\S8.1, properties (1)--(2)]{Zhang}, with the usual elliptic
Kummer conditions away from the support. Assume also that
\[
 \dim_{\mathbb F_7}\operatorname{Sel}_7(E/K)^+=4.
\]
For the support $n_*$ assume that $\kappa_{n_*}$ lies in the
plus eigenspace and that $\kappa_{n_*/\ell}=0$ for each
$\ell\mid n_*$. In particular the latter assumption holds if the
system's minimal nonzero support has three prime factors.

To see precisely why the \emph{entire} localization vanishes, let
$v$ be the unique place of $K$ over $\ell\mid n_*$. The finite
(unramified) and transverse local conditions are complementary:
\[
 H^1(K_v,E[7])=H^1_{\rm fin}(K_v,E[7])
                   \oplus H^1_{\rm tr}(K_v,E[7]).
\]
By property (1), $\operatorname{loc}_v\kappa_{n_*}$ already belongs
to the transverse summand. Its projection to the singular quotient
$H^1/H^1_{\rm fin}$ determines it uniquely there. Under this
identification, equation (8.1) of \cite{Zhang} gives
\[
 \operatorname{loc}_v\kappa_{n_*}
 =\psi_\ell(\operatorname{loc}_v\kappa_{n_*/\ell})=0,
 \qquad \ell\mid n_*.
\]
Thus the class satisfies the ordinary Selmer condition at every
place, including all three support places. Its plus Selmer space
is spanned by the four rational point classes: their mod-seven
independence persists after restriction to $K$ by
restriction--corestriction, and the assumed dimension is four.

Applying the three rows of the localization matrix gives
\[
 \boxed{
 \kappa_{n_*}
   =\lambda\,\delta_7(B),
 \qquad\lambda\in\mathbb F_7.
 }
\]

When the minimal nonzero support has size three, at least one
three-prime support gives a nonzero class. This leaves the value
of $\lambda$ at $n_*$ open.

\section{Two exact finite-field detectors}

The prime 167 is inert in $K$, and the
$K$-prime $(167)$ is principal and coprime
to $n_*$. It therefore splits completely
in $L/K$.

One computes
\[
 \#E(\mathbb F_{167^2})
   =168^2-14^2
   =28028=49\cdot572.
\]
Its 7-primary subgroup is $E[7]$.

On the short model,
\[
 \overline B_{167}=(152,63),\qquad
 [572]\overline B_{167}=(166,31)\ne O.
\]
The required scalar is consequently
characterized by
\[
 [572]\overline z_{167}
    =[\lambda](166,31),
\]
provided $\overline z_{167}$ is the reduction
of the correctly normalized derived point.

A smaller detector is available at 19.
It is also inert in $K$, with principal
$K$-prime $(19)$, and splits completely in
$L/K$. It need not be Kolyvagin-admissible
because it is used only for detection.

Here
\[
 \#E(\mathbb F_{19})=28,\qquad
 \#E(\mathbb F_{19^2})=336=7\cdot48.
\]
On the short model,
\[
 \overline B_{19}=(1,2),\quad
 [4]\overline B_{19}=(2,13),\quad
 V_{19}:=[48]\overline B_{19}=(3,2)\ne O.
\]
The group $E(\mathbb F_{19^2})[7]$ is the
cyclic group generated by $V_{19}$.

Define
\[
 L_{19}:E(\mathbb F_{19^2})/7
           \longrightarrow\mathbb F_7,
 \qquad
 [48]Q=[L_{19}(Q)]V_{19}.
\]

Then the exact unevaluated scalar is
\[
 \boxed{
 \lambda=
 \sum_{r_1,r_2,r_3=1}^{6}
 r_1r_2r_3\,
 L_{19}\!\left(
 \operatorname{red}_{19}
 (g_1^{r_1}g_2^{r_2}g_3^{r_3}Y_H)
 \right).
 }
\]

Equivalently,
\[
 \kappa_{n_*}\ne0
 \quad\Longleftrightarrow\quad
 [48]\overline z_{19}\ne O.
\]

\section{Finite reductions in a nodal comparison}

For the same four points one obtains a second independence
certificate at $3$. With the coordinate of the first nonzero local
point normalized to one, at $\ell=11,17,23,29$ the rows in
$E(\mathbb F_\ell)/3$ are
\[
 \begin{pmatrix}
 1&0&0&1\\0&1&0&1\\0&1&2&0\\1&0&0&2
 \end{pmatrix}.
\]
The group orders are $18,24,30,36$, respectively. Each prime is
$2\bmod3$, so again the endomorphism $[\#E(\mathbb F_\ell)/3]$
detects the one-dimensional quotient. The matrix is invertible.

The mod-three Galois image for this curve is full \cite{LMFDB}.
It also follows by factoring its three-division polynomial
\[
 3X^4-3X^3-474X^2+3468X-7108
\]
modulo $5$ and $7$, obtaining cycle types $(4)$ and $(1)(3)$.
The only subgroup of $S_4$ containing both types is $S_4$.
A lift of a projective four-cycle in $\operatorname{GL}_2(\mathbb F_3)$
has fourth power $-I$ (use the nonsplit multiplicative group
$\mathbb F_9^\times$), so the kernel also contains the scalar $-I$.

For $1\le k\le4$ set
\[
 D_k=\mathbb Q(E[3],3^{-1}P_1,\ldots,3^{-1}P_k).
\]
Then $[D_k:\mathbb Q]=48\cdot3^{2k}$. Indeed $-I$ acts by $-1$
on $E[3]$ and trivially by conjugation on its Galois group, implying
$H^1(\operatorname{GL}_2(\mathbb F_3),E[3])=0$. Restriction of the
independent Kummer classes to the three-division field remains
independent. The translation image is a submodule of $E[3]^k$. If
it were proper, a nonzero equivariant map from $E[3]^k$ to $E[3]$
would kill it, because this is a sum of simple modules with
endomorphisms $\mathbb F_3$. This would give a relation among the
restricted classes. Thus, by Chebotarev, good primes where $E[3]$
is rational and all $k$ points are three-divisible in the residue
group have density $1/(48\cdot3^{2k})$.

A distinction between finite group size and Kummer data can be seen
in the generalized Jacobian $J$ of the curve obtained by gluing
$P_i$ to $2P_i$ for $i=1,2,3$. At $p=373$ these six branches remain
distinct. The extension by $\mathbb G_m^3$ has classes $-P_i$ in
$\operatorname{Pic}^0(E)$. The finite-field connecting map
\[
 E(\mathbb F_p)[3]\longrightarrow
 \bigl(\mathbb F_p^\times/(\mathbb F_p^\times)^3\bigr)^3
\]
is given by the finite Tate pairings with these classes. This is the
definition of the Tate pairing via lifting a torsion point to the
extension of the Jacobian by the line bundle. The pairing with
$E(\mathbb F_p)/3$ is perfect since $p\equiv1\bmod3$. Enumeration
on the short model gives $\#E(\mathbb F_{373})=405$ and
$\#(3E(\mathbb F_{373}))=45$. Modulo $3E(\mathbb F_{373})$, the first
two points form a basis and $P_3$ has coordinates $(2,2)$. These
facts can be reproduced by tripling all $405$ points and comparing
$P_i-uP_1-vP_2$ with this set for $0\le u,v\le2$. The map therefore
has rank two. It follows that $\#J(\mathbb F_{373})[3]=27$, whereas
the split extension has $243$ points killed by three. Both groups
have total $3$-primary order $3^7$. Their $3^m$-torsion
cardinalities agree for $m\ge7$.

More generally, for $k$ split nodes the order
$(p-1)^k\#E(\mathbb F_p)$ is independent of the gluing points. If
$\ell^m>(Y-1)^k(\sqrt Y+1)^2$, the $\ell^m$-torsion cardinalities
of the extension and its split counterpart agree simultaneously
at all good primes $p\le Y$ with these split nodes. This follows by
bounding the order of each group and then taking its full
$\ell$-primary subgroup.

\section{Arithmetic scope}
The detector at $19$ reduces a class on the specified Kummer line
to one scalar. Its evaluation requires the actual conductor-$n_*$
CM points and their integral elliptic quotient, including the
chosen generators and orientations. The sign or nonvanishing of this
scalar has not been established. The conditional formula and the
finite local certificates above yield no determination of the
central second derivative.
\begin{thebibliography}{9}

\bibitem{LMFDB}
LMFDB Collaboration, \emph{Elliptic curve 234446.a1}, including its
residual Galois images. \url{https://www.lmfdb.org/EllipticCurve/Q/234446/a/1}.
\bibitem{Zhang}
W. Zhang,
\emph{Selmer groups and the indivisibility
of Heegner points},
Cambridge Journal of Mathematics 2 (2014),
191--253, especially Sections 3 and 8.
\url{https://math.mit.edu/~wz2113/math/online/Kconj.pdf}

\end{thebibliography}

\end{document}
